\section{The Judge Who Was Judged}\label{sec:judge}

``For neither does the Father judge any man: but hath given all judgment to the Son'' (John 5:22). And a few verses later: ``he hath given him power to do judgment, because he is the Son of man'' (5:27). The Judge before whom every soul will stand at death, and before whom the whole race will be gathered at the end, is the Son, the Word made flesh, who has a human face and bears in his hands and side the marks of his Passion. The Church sings it in the \latin{Te Deum}: \latin{Judex crederis esse venturus}: thou art believed to be the Judge who is to come.\endnote{\latin{Te Deum}, Latin as printed in John Henry Blunt (ed.), \emph{The Annotated Book of Common Prayer} (London, 1866), p.~12.} And at once she adds a petition, which Alphonsus takes into his own prayer: \latin{Te ergo quaesumus tuis famulis subveni, quos pretioso sanguine redemisti}---``We therefore beseech thee, assist thy servants, whom thou hast redeemed with thy precious blood.''\endnote{Alphonsus Liguori, \emph{Preparation for Death}, Consideration~XVI, Affections and Prayers (Duffy, 1843), p.~204.}

\subsection{Why the Son of Man judges}

Aquinas asks why judgment belongs to Christ not only as God but also as man. The ground he gives first is headship: Christ even in his human nature is head of the whole Church, and the power to judge belongs to the head. He then gives three reasons why it is fitting. The second is that at the last day the bodies of the dead will be raised through the Son of man; the third, that those who are to be judged should see their Judge. The first gives a sinner ground for confidence: judiciary power belongs to Christ in his human nature ``because of His likeness and kinship with men; for, as God works through intermediary causes, as being closer to the effects, so He judges men through the Man Christ, that His judgment may be sweeter to men.''\endnote{\emph{Summa theologiae} III, q.~59, a.~2, corpus (English Dominican translation).} He then cites the Epistle to the Hebrews: ``For we have not a high priest who cannot have compassion on our infirmities: but one tempted in all things like as we are, without sin. Let us go therefore with confidence to the throne of grace'' (Heb 4:15--16).

In the next article Aquinas adds that Christ received this power not only by right of his divinity and his headship but also by merit, ``so that, in accordance with the Divine justice, He should be judge who fought for God's justice, and conquered, and was unjustly condemned.''\endnote{\emph{Summa theologiae} III, q.~59, a.~3, corpus.} The Catechism says the same: ``Full right to pass definitive judgement on the works and hearts of men belongs to him as redeemer of the world. He `acquired' this right by his cross.''\endnote{\emph{Catechism of the Catholic Church} 679, official English.} The Judge of the living and the dead is the one who stood before Pilate and was condemned.

The \latin{Dies irae} calls this Judge the \latin{Rex tremendae majestatis}, and then appeals to him by what he has already done:

\begin{quote}
\latin{Quaerens me, sedisti lassus:\\
\quad Redemisti, crucem passus:\\
\quad Tantus labor non sit cassus.}
\end{quote}

\noindent The singer reminds Christ that he sat down weary in seeking him and redeemed him by suffering the cross, and begs that so great a labour may not be in vain. The words echo John's account of Jesus ``wearied with his journey,'' sitting on Jacob's well where the Samaritan woman came to draw water (John 4:6--7). The Judge to whom the singer appeals is, in this strophe, the one who came looking for him, like the shepherd who went after the lost sheep.

\subsection{The least of the brethren}

When Christ describes his judgment in Matthew~25, the matter of it is plain: whether the nations fed him when he was hungry and visited him when he was sick. John Chrysostom, preaching on the passage, remarks on how little is asked:

\begin{quote}
And mark how easy are His injunctions. He said not, ``I was in prison, and ye set me free; I was sick, and ye raised me up again;'' but, ``ye visited me,'' and, ``ye came unto me.'' And neither in hunger is the thing commanded grievous. For no costly table did He seek, but what is needful only, and His necessary food, and He sought in a suppliant's garb \ldots\endnote{John Chrysostom, \emph{Homilies on Matthew}, Homily~79, 1, trans.\ George Prevost, rev.\ M.~B. Riddle, \emph{Nicene and Post-Nicene Fathers}, first series, vol.~10 (New York, 1888).}
\end{quote}

Chrysostom also says who the least of the brethren are: ``Why, for this reason they are brethren, because they are lowly, because they are poor, because they are outcast. For such doth He most invite to brotherhood, the unknown, the contemptible.'' He refuses to restrict the word to monks or to the holy: ``every believer; though he be a secular person, yet if he be hungry, and famishing, and naked, and a stranger, His will is he should have the benefit of all this care. For baptism renders a man a brother, and the partaking of the divine mysteries.''\endnote{Chrysostom, \emph{Homilies on Matthew}, Homily~79, 1.} The Catechism, expounding the last judgment, quotes Augustine, who has Christ say: ``I placed my poor little ones on earth for you. I as their head was seated in heaven at the right hand of my Father---but on earth my members were suffering, my members on earth were in need. If you gave anything to my members, what you gave would reach their Head.''\endnote{Augustine, \emph{Sermo} 18, 4 (PL 38, 130--131), as quoted at \emph{Catechism of the Catholic Church} 1039, official English.}

In the scene both the blessed and the cursed are surprised. ``Lord, when did we see thee hungry and fed thee: thirsty and gave thee drink?'' ask the just; ``Lord, when did we see thee hungry \ldots\ and did not minister to thee?'' ask the others (Matt 25:37, 44). Neither group knew that it was serving or refusing Christ. The just do not recognize their own record, because charity does not keep accounts of itself.

The two sentences also differ. To the blessed the King says, ``Come, ye blessed of my Father, possess you the kingdom prepared for you from the foundation of the world.'' To the others, ``Depart from me, you cursed, into everlasting fire, which was prepared for the devil and his angels.'' Chrysostom reads them word by word. The cursed are ``no longer of the Father; for not He laid the curse upon them, but their own works''; and the fire was not prepared for them:

\begin{quote}
\ifdefined\TriptychPrintEdition
\looseness=1
\else
\fi
For concerning the kingdom indeed, when He had said, ``Come, inherit the kingdom,'' He added, ``prepared for you before the foundation of the world;'' but concerning the fire, no longer so, but, ``prepared for the devil.'' I, saith He, prepared the kingdom for you, but the fire no more for you, but ``for the devil and his angels;'' but since ye cast yourselves therein, impute it to yourselves.\endnote{Chrysostom, \emph{Homilies on Matthew}, Homily~79, 2.}\par
\end{quote}

It follows that God's first intention for every soul is the kingdom, and that the Judge who pronounces the sentence is the one who prepared it. Benedict~XVI ends his reflection on the last judgment with the counsel that grace ``allows us all to hope, and to go trustfully to meet the Judge whom we know as our `advocate', or \emph{parakletos}.''\endnote{\emph{Spe salvi} 47, official English.} So the Apostle John had written: ``if any man sin, we have an advocate with the Father, Jesus Christ the just'' (1~John 2:1).
